\section{Introduction}\label{sec:introduction}

At the critical point of the planar Ising model, a boundary between plus
and minus spins has a conformally invariant scaling limit. Independent
random variations of the bond strengths break the spatial symmetries in
every typical realization. The question considered here is whether, when
the variations are sufficiently weak, the interface in a fixed environment
still approaches the same deterministic probability law.

We prove this for symmetric two-valued ferromagnetic bonds. The disorder
strength is held fixed while the mesh tends to zero. The limiting object
is the full oriented chordal curve, and the temperature is specified by
the onset of spontaneous magnetization.

\subsection{The random ferromagnet and its critical point}

Let $E^\bullet$ be the set of unordered nearest-neighbor edges of $\Z^2$,
each counted once. On a probability space
$(\Omega_{\rm env},\Prob_{\rm env})$, let $(\xi_e)_{e\in E^\bullet}$
be independent with
\[
 \Prob_{\rm env}(\xi_e=1)=\Prob_{\rm env}(\xi_e=-1)=\tfrac12.
\]
Fix $0<\epsilon<1$ and set
\begin{equation}\label{intro:bonds}
 J_e(\omega)=1+\epsilon\xi_e(\omega).
\end{equation}
The same infinite environment is used for all finite volumes. There is
no magnetic field. For $\Lambda_m=\{-m,\ldots,m\}^2$, extend spins by
$+1$ outside $\Lambda_m$ and define
\[
 H^+_{m,\omega}(\sigma)
   =-\sum_{e=\{x,y\}\in E^\bullet:\ e\cap\Lambda_m\ne\varnothing}
                    J_e(\omega)\sigma_x\sigma_y.
\]
Write $\langle\cdot\rangle^+_{m,\beta,\omega}$ for the normalized
Gibbs expectation with weight $e^{-\beta H^+_{m,\omega}}$. Put
\begin{equation}\label{intro:critical-definition}
 m_\epsilon(\beta)
 =\E_{\rm env}\!\left[\lim_{m\to\infty}
                     \langle\sigma_0\rangle^+_{m,\beta,\omega}\right],
 \qquad
 \beta_c(\epsilon)=\inf\{\beta\ge0:m_\epsilon(\beta)>0\}.
\end{equation}
Ferromagnetic monotonicity gives the limit and, with
$K_0=\tfrac12\log(1+\sqrt2)$,
\begin{equation}\label{intro:critical-bounds}
 0<\frac{K_0}{1+\epsilon}\le\beta_c(\epsilon)
           \le\frac{K_0}{1-\epsilon}<\infty.
\end{equation}
The Gibbs normalization in this definition, and in every finite-volume
law below, is taken separately for each environment.

\subsection{Domains and interfaces}

The dual lattice has vertex set
$F=\Z^2+(\tfrac12,\tfrac12)$ and nearest-neighbor edge set $E$.
For $e\in E$, write $e^\bullet$ for the crossing primal edge.
If $V\subset\Z^2$ is finite, set
\[
 Q_V=\bigcup_{v\in V}\bigl(v+[-\tfrac12,\tfrac12]^2\bigr),
 \qquad D_V=\operatorname{int}Q_V.
\]
We require $Q_V$ to be the closure of a Jordan domain. Let $E_V$ be
the dual edges crossing primal bonds with at least one endpoint in $V$.
Choose distinct marks $f_{\rm in},f_{\rm out}\in F\cap\partial D_V$.
At each mark choose an exterior stub: a dual edge outside $E_V$ joining
that mark to a vertex outside $Q_V$. Denote these stubs by
$e_{\rm in}$ and $e_{\rm out}$.

For an edge set $A$, let $d_f(A)$ be its degree at $f$. The permitted
interior contour sets are
\begin{equation}\label{intro:contour-space}
 \mathcal P_V=\left\{P\subset E_V:
 d_f(P)+\1_{\{f=f_{\rm in}\}}+\1_{\{f=f_{\rm out}\}}
                \text{ is even for every }f\in F\right\}.
\end{equation}
Require $\mathcal P_V\ne\varnothing$ and the existence of a finite
completion $P_o\subset E\setminus E_V$ containing both stubs, with
the same two odd-degree marks and even degree elsewhere. These choices
will be called \emph{admissible lattice data}. The completion is fixed
deterministically. For $P\in\mathcal P_V$, the union $P\cup P_o$
is even and determines spins equal to plus far away, changing sign
across its edges. The resulting exterior spins are independent of $P$
and give opposite boundary signs on the two marked arcs.

For a fixed environment, at inverse temperature $\beta$, use the law
\begin{equation}\label{intro:quenched-contours}
 \mu_{V,\beta,\epsilon,\omega}(P)
 =\frac1{Z_{V,\beta,\epsilon,\omega}}
       \exp\left(-2\beta\sum_{e\in P}J_{e^\bullet}(\omega)\right),
 \qquad P\in\mathcal P_V,
\end{equation}
where the denominator is the sum of the numerator over
$\mathcal P_V$. In particular, the weights include bonds crossing
the boundary of $V$.

\begin{lemma}[Contour and spin laws]\label{intro:contour-spin}
The law \eqref{intro:quenched-contours} is the contour image of the
nearest-neighbor Ising law on $V$ with the exterior spins determined by
$P_o$. It is independent of the choice of completion $P_o$.
\end{lemma}

\begin{proof}
An even finite dual edge set determines a unique spin configuration
equal to plus far away: the sign at a primal vertex is the parity of
crossings of any path from that vertex to infinity. Evenness makes this
parity independent of the path. Restricting the contour to $E_V$ gives
a bijection between the compatible spin configurations on $V$ and
$\mathcal P_V$. Across a primal bond, the spin product is $-1$
exactly when its crossing dual edge belongs to the contour. Thus
\[
 \exp\left(\beta\sum_{\{x,y\}:\{x,y\}\cap V\ne\varnothing}
                      J_{\{x,y\}}\sigma_x\sigma_y\right)
 =\exp\left(\beta\sum_{e\in E_V}J_{e^\bullet}\right)
       \exp\left(-2\beta\sum_{e\in P}J_{e^\bullet}\right).
\]
The first factor is constant over $P$. Every exterior-only energy is
constant as well. Normalization therefore gives
\eqref{intro:quenched-contours}, whose formula does not involve $P_o$.
\end{proof}

Add to $P$ the half of each stub from its mark to its midpoint. At
degree two join the two occupied edges. At degree four, pair north
with east and south with west, also at marked vertices. Draw these
pairings disjointly in disks of radius less than $1/10$ about the dual
vertices, using any fixed local drawing. The resolved component from
the incoming stub midpoint to the outgoing midpoint is a simple
oriented curve, denoted $\gamma(P)$.

\subsection{The limiting law}

For continuous oriented curves in $\C$, represented on $[0,1]$, let
\begin{equation}\label{intro:curve-distance}
 d_{\rm curv}(\gamma,\eta)
 =\inf_{\alpha,\rho}\sup_{t\in[0,1]}
             |\gamma(\alpha(t))-\eta(\rho(t))|,
 \qquad \bar d=\min\{1,d_{\rm curv}\},
\end{equation}
where $\alpha,\rho$ run over increasing homeomorphisms of $[0,1]$;
curves at distance zero are identified. For probability measures
$\nu,\nu'$ on this curve space define
\begin{equation}\label{intro:BL-distance}
 d_{\rm BL}(\nu,\nu')
 =\sup_{\substack{|F|\le1\\
           |F(\gamma)-F(\eta)|\le\bar d(\gamma,\eta)}}
           \left|\int F\,\dd\nu-\int F\,\dd\nu'\right|.
\end{equation}

Let $D$ be a bounded Jordan domain and $a,b$ distinct boundary points.
Chordal $\SLE_3(D;a,b)$ is the law obtained by conformally mapping the
upper-half-plane Loewner trace
\[
 \partial_t g_t(z)=\frac2{g_t(z)-\sqrt3 B_t},\qquad g_0(z)=z,
\]
to $D$, with $0$ mapped to $a$ and $\infty$ to $b$; $B_t$ is standard
real Brownian motion. We retain the oriented curve modulo
reparametrization.

\begin{theorem}[Quenched interface universality]\label{thm:main}
There is $\epsilon_0>0$ such that the following holds for every fixed
$\epsilon\in(0,\epsilon_0)$. Let $D$ be a bounded Jordan domain and
$a,b\in\partial D$ distinct. Take any deterministic sequence of admissible
lattice data with meshes $\delta_n\downarrow0$ and boundary
homeomorphisms
\[
 \phi_n:\partial D\longrightarrow\partial(\delta_n D_{V_n}),
 \qquad \sup_{z\in\partial D}|\phi_n(z)-z|\longrightarrow0,
\]
such that $\phi_n(a)=\delta_n f_{{\rm in},n}$ and
$\phi_n(b)=\delta_n f_{{\rm out},n}$. Let $Q_{n,\omega}$ be the law
of $\delta_n\gamma(P)$ under
$\mu_{V_n,\beta_c(\epsilon),\epsilon,\omega}$. Then for every $\eta>0$,
\begin{equation}\label{intro:main-convergence}
 \Prob_{\rm env}\!\left[
 d_{\rm BL}\bigl(Q_{n,\omega},\SLE_3(D;a,b)\bigr)>\eta\right]
                   \longrightarrow0.
\end{equation}
\end{theorem}

The domain approximations are independent of the environment, and the
convergence uses their original Euclidean embedding. The strength
$\epsilon$ is fixed throughout the limit. The proof also identifies
the critical point: for sufficiently small $\epsilon$,
\begin{equation}\label{intro:selfdual-beta}
 \bigl(e^{2\beta_c(\epsilon)(1+\epsilon)}-1\bigr)
 \bigl(e^{2\beta_c(\epsilon)(1-\epsilon)}-1\bigr)=2.
\end{equation}
This identity is proved through magnetization and random-cluster
crossing estimates in Section~\ref{sec:criticality}.

\subsection{History and related work}

Harris's criterion gives a basic test for the relevance of quenched
disorder~\cite{Harris}. It leaves the
marginal case of the two-dimensional Ising model undecided. The
renormalization-group analysis of Dotsenko and Dotsenko
\cite{DD,DDReview} predicted a double-logarithmic specific-heat
singularity. Subsequent work of Shalaev, Shankar, and Ludwig
\cite{Shalaev,Shankar,Ludwig} developed the picture of unchanged leading
critical exponents together with logarithmic corrections, revising the
earlier prediction for spin correlations. This is the physical setting
of marginal irrelevance: weak bond disorder decreases under coarse graining, but
only slowly enough that its cumulative effect can remain visible in
logarithmic factors. These calculations concern thermodynamic quantities
and correlations; a quenched law for an entire interface requires
additional geometric control.

For the homogeneous critical model, the conformal description is
rigorous. Schramm introduced stochastic Loewner evolution, which now describes
many planar interface limits~\cite{Schramm}. Smirnov proved conformal
convergence of the FK-Ising fermionic observable~\cite{Smirnov}, and
Chelkak and Smirnov developed fermionic-observable universality on
critical isoradial graphs~\cite{ChelkakSmirnov}. Chelkak,
Duminil-Copin, Hongler, Kemppainen, and Smirnov then proved convergence
of the critical square-lattice Dobrushin spin interface to chordal
$\SLE_3$~\cite[Theorem~1 and Equation~(1)]{CDHKS}. Their conclusion is
convergence of curves in uniform distance modulo orientation-preserving
reparametrization, the topology used here.

Two other directions provide relevant comparisons. For weak, translation- and rotation-invariant
finite-range pair perturbations, Giuliani, Greenblatt, and
Mastropietro proved universality of infinite-plane energy-correlation
scaling limits~\cite[Theorems~1.1--1.2]{GGM}. Greenblatt and Peltola
constructed a Grassmann martingale observable for non-integrable spin
interfaces and formulated asymptotic conditions that would imply an
$\SLE_3$ limit~\cite{GP}. The companion~\cite{S} supplies the
deterministic transfer and interface comparison results used in this
paper. For random couplings, Mahfouf proved quenched FK-interface
convergence to $\SLE_{16/3}$ in a regime in which the disorder strength
vanishes with the mesh~\cite[Theorem~1.6]{Mahfouf}; Av\'erous and
Mahfouf proved quenched crossing estimates for near-critical random
bond FK-percolation in related regimes~\cite[Theorem~1.6]{AM}.
The disorder distribution in the present theorem is fixed as the mesh
tends to zero, and the conclusion concerns spin rather than FK interfaces.

Rigorous results on disordered phase transitions also distinguish the
location of a transition from its scaling behavior. Aizenman and Wehr
proved that quenched bond randomness rounds energy discontinuities in
two-dimensional ferromagnets at positive temperature
\cite[Theorem~4.1]{AW}. Chayes and Shtengel proved zero magnetization,
uniqueness of the limiting disorder-averaged state, and divergent
averaged susceptibility at self-dual iid
ferromagnetic coupling laws~\cite[Theorems~$1'$--$2'$]{ChayesShtengel}.
Their result does not by itself exclude a critical interval containing
the self-dual point. In this paper the identification with the physical
magnetization threshold is proved separately from the construction of
the critical scale trajectory.

The proof also uses established geometric and probabilistic methods.
Kramers--Wannier duality~\cite{KW}, including its four-sector torus
form~\cite[Equations~(4)--(5)]{BS}, identifies the symmetric tuning.
The random-cluster representation of Fortuin and Kasteleyn~\cite{FK},
the Edwards--Sokal coupling~\cite{ES}, and positive association
\cite{FKG} connect spin and crossing events. The decision-tree
inequality of O'Donnell, Saks, Schramm, and Servedio~\cite{OSSS} and
its extension to monotonic measures by Duminil-Copin, Raoufi, and
Tassion~\cite[Theorem~1.1]{DCRT} supply the sharpness argument at fixed environment, after which we average
over the disorder. Finally, the use of crossing estimates to
control curves has precedents in Aizenman--Burchard~\cite{AB} and
Kemppainen--Smirnov~\cite{KS}; Sheffield--Sun~\cite{SS} explain why
convergence of Loewner driving functions alone does not supply the full
curve topology. Here the needed geometric comparison and exclusion of
macroscopic reverse traversals are proved within the argument.


\subsection{Proof strategy and the new estimates}

Our starting point is the pure-model analysis and deterministic geometry
of the companion paper~\cite{S}. We use its pure conditional-law
comparisons, transfer estimates for spatially symmetric observables,
crossing and passage bounds, and geometric passage from boundary barriers
to curves. Section~\ref{sec:pure} derives the additional transfer and
separated-product estimates needed for observables without spatial
symmetry. A typical random realization has no such symmetry, so the
subsequent scale analysis must control a random family of local
interactions uniformly in the exterior field data.

The scale transformation keeps the interactions local by adding
auxiliary fair coins. These coins belong to the reference Gibbs field;
they are distinct from the random environment. After conditioning on
them, a separated collection of squares can be integrated exactly.
The resulting interactions depend only on bounded neighborhoods of
the enlarged scale. A separate operation treats isolated concentrations
of large interactions, so estimates use environmental moments rather
than a supremum over environments.

The relevant local coordinate is the coefficient of the pure energy
field. We construct a linear functional $Q$ on bounded even local
observables that is exactly preserved by transfer through a pure
annulus. A covariance sum of these charges measures the leading
disorder strength. Its variance receives a negative quadratic correction
under one scale step. The sign comes from the normalization of a joint
insertion of two separated interactions. A finite polynomial correction
to the covariance sum makes this calculation compatible with the local
transformation and its allocation of interactions to larger cells.

The resulting disorder decays only logarithmically with length scale.
For crossings, this is sufficient because a change in a small square
can alter a macroscopic test only when several paths reach away from
that square. Pure passage bounds supply a positive power of the ratio
of the square size to the test size. Summing these geometrically
weighted errors makes the full comparison tend to zero. We first
take a supremum over exterior spin data; the resulting local functions
of the environment can then be multiplied for disjoint ring
neighborhoods. This order of operations gives the conditional estimates
needed both for criticality and for boundary comparisons.

Finally, the target interface is coupled to pure comparison chords whose
law is deterministic even after conditioning on the environment. The
expected distance in these couplings bounds the expected
bounded-Lipschitz distance between the quenched laws. Boundary barriers
and an interior multiple-passage estimate turn this comparison into
convergence of the full oriented curve.

The exact local transformation, its treatment of large interactions,
and the polynomial variance coordinate provide reusable methods for
random perturbations with a marginal energy direction. Their use here
requires the stated pure transfer and geometric estimates. All scale
parameters and the allowed disorder interval are fixed before the
target domain is chosen.

\paragraph{Organization.}
Section~\ref{sec:pure} develops the pure energy charge and its product
estimates. Sections~\ref{sec:local-map} and~\ref{sec:moments} construct
the local transformation and prove its moment bounds.
Sections~\ref{sec:variance} and~\ref{sec:trajectory} establish the
variance recursion and tune a trajectory with decaying disorder.
Section~\ref{sec:crossings} transfers those bounds to crossings and
passages. Section~\ref{sec:criticality} identifies the physical critical
point. Section~\ref{sec:curves} proves Theorem~\ref{thm:main}.
